% !TeX program = xelatex
% !LW recipe = latexmk (xelatex)
\synctex=1
\RequirePackage{fix-cm}
\documentclass[12pt,a4paper,landscape]{article}
\usepackage[margin=0mm]{geometry}
\usepackage{fontspec}
\setmainfont{DejaVuSans}[
  Extension=.ttf,
  UprightFont=*,
  BoldFont=*-Bold,
  ItalicFont=*-Oblique,
  BoldItalicFont=*-BoldOblique]
\usepackage[russian]{babel}
\usepackage{tikz}
\usetikzlibrary{arrows.meta,calc}
\usepackage{qrcode}
\usepackage{hyperref}
\hypersetup{hidelinks,pdftitle={Сортировщик для мини-склада},
  pdfauthor={Малышев Павел; Войлов Антон; Янковец Александр},
  pdfsubject={Лифлет проекта ППП}}
\pagestyle{empty}
\setlength{\parindent}{0pt}
\setlength{\parskip}{0pt}
\hyphenpenalty=10000
\exhyphenpenalty=10000
\definecolor{ink}{gray}{0.08}
\definecolor{muted}{gray}{0.30}
\definecolor{rulegray}{gray}{0.68}
\newcommand{\bodyFont}{\fontsize{12bp}{15bp}\selectfont}
\newcommand{\sectionLabel}[2]{{\fontsize{13bp}{16bp}\selectfont\bfseries #1\enspace #2}}
\newcommand{\subLabel}[1]{{\bfseries #1}\par\vspace{2mm}}
\newcommand{\commandText}[1]{\mbox{\textbf{#1}}}
% All coordinates below are millimetres from the top-left of the page.
\newcommand{\putBlock}[4]{%
  \node[anchor=north west,inner sep=0pt,outer sep=0pt,
        text width=#3mm,align=left,font=\bodyFont] at (#1,-#2)
        {\raggedright #4};}
\newcommand{\pageFrame}[2]{%
  \putBlock{12}{9}{210}{\textbf{SORTING WAREHOUSE}\enspace /\enspace ППП}
  \draw[line width=.6pt] (12,-18)--(285,-18);
  \draw[rulegray,line width=.4pt] (12,-198)--(285,-198);
  \putBlock{12}{201}{240}{\color{muted}#1}
  \putBlock{266}{201}{20}{\hfill #2\,/\,2}}
\newcommand{\teamMember}[4]{%
  \putBlock{#1}{162}{54}{\textbf{#2}\\#3\par\vspace{2mm}#4}}
\newcommand{\stepBlock}[3]{%
  \putBlock{193}{#1}{8}{\textbf{#2}}
  \putBlock{203}{#1}{82}{#3}}
\begin{document}
\bodyFont\color{ink}
\null
\begin{tikzpicture}[remember picture,overlay,x=1mm,y=1mm]
\begin{scope}[shift={(current page.north west)}]
  \pageFrame{Учебный прототип автоматизированного склада}{1}
  \putBlock{12}{25}{108}{\sectionLabel{01}{О проекте}}
  \putBlock{12}{35}{110}{{\fontsize{26bp}{29bp}\selectfont\bfseries
    Сортировщик\\для мини-склада}}
  \putBlock{12}{61}{101}{Учебный прототип принимает предметы с конвейера
    и распределяет их по девяти ячейкам. Камера определяет цвет, весы измеряют
    массу, а головка переносит предмет в выбранный ряд.}
  \putBlock{12}{102}{108}{\sectionLabel{04}{Инженерная задача}}
  \putBlock{12}{112}{102}{Согласовать распознавание, взвешивание и перемещение
    в одном цикле. Главный вызов — снизить влияние вибраций на весы и подать
    предмет точно в приёмную головку.}

  \putBlock{127}{25}{158}{\sectionLabel{06}{Устройство в сборе}}
  \node[anchor=north west,inner sep=0pt,outer sep=0pt] at (127,-34)
    {% Photo-derived axonometric sketch; units in mm, no dimensional claims.
% Requires TikZ only. Labels: 1 conveyor, 2 camera, 3 XY head,
% 4 storage cells, 5 control electronics. Bounding box: 158 x 88 mm.
\begin{tikzpicture}[x=1mm,y=1mm,line cap=round,line join=round,
  line width=.45pt,draw=black!80]
\path[use as bounding box] (0,0) rectangle (158,88);
\begin{scope}[shift={(9,20)},x={(.85mm,-.16mm)},
  y={(.46mm,.63mm)},z={(0mm,.95mm)}]
  % Frame and corner supports, far edges first.
  \draw[fill=white] (0,58,-2)--(80,58,-2)--(80,62,-2)--(0,62,-2)--cycle;
  \draw (0,58,-2)--(0,58,2)--(80,58,2)--(80,58,-2);
  \draw[fill=white] (78,0,-2)--(82,0,-2)--(82,60,-2)--(78,60,-2)--cycle;
  \draw (82,0,-2)--(82,0,2)--(82,60,2)--(82,60,-2);
  \draw[fill=white] (0,0,-2)--(4,0,-2)--(4,60,-2)--(0,60,-2)--cycle;
  \draw (0,0,-2)--(0,0,2)--(0,60,2)--(0,60,-2);
  \foreach \frameX in {0,76}{
    \foreach \frameY in {0,56}{
      \draw[fill=white] (\frameX,\frameY,2)--++(5,0,0)--++(0,5,0)--++(-5,0,0)--cycle;
      \draw (\frameX,\frameY,2)--++(0,0,5)--++(5,0,0)--++(0,5,0)--++(0,0,-5);
      \draw (\frameX,\frameY,7)--++(0,5,0)--++(5,0,0);
    }
  }
  % Nine-cell tray with a visible open top and low outside walls.
  \draw[fill=white,line width=.65pt]
    (30,10,0)--(72,10,0)--(72,52,0)--(30,52,0)--cycle;
  \draw[fill=white] (30,10,0)--(72,10,0)--(72,10,7)--(30,10,7)--cycle;
  \draw[fill=white] (72,10,0)--(72,52,0)--(72,52,7)--(72,10,7)--cycle;
  \draw[line width=.65pt] (30,10,7)--(72,10,7)--(72,52,7)--(30,52,7)--cycle;
  \draw (30,10,0)--(30,52,0)--(30,52,7);
  \foreach \cellX in {44,58}{
    \draw[line width=.65pt] (\cellX,10,7)--(\cellX,52,7);
    \draw[black!35] (\cellX,10,7)--(\cellX,10,1);
  }
  \foreach \cellY in {24,38}{
    \draw[line width=.65pt] (30,\cellY,7)--(72,\cellY,7);
    \draw[black!35] (72,\cellY,7)--(72,\cellY,1);
  }
  % Far guide rod and its adjacent belt.
  \draw[line width=1.8pt,black!60] (3,60,8)--(79,60,8);
  \draw[line width=.45pt,white] (3,60,8)--(79,60,8);
  \draw[line width=.9pt] (3,56,9)--(79,56,9);
  % Electronics platform: power module and two driver boards.
  \draw[fill=white] (2,14,5)--(18,14,5)--(18,53,5)--(2,53,5)--cycle;
  \draw (2,14,5)--(2,14,1)--(18,14,1)--(18,53,1)--(18,53,5);
  \foreach \boardY in {17,29,41}{
    \draw[fill=white] (4,\boardY,6)--++(12,0,0)--++(0,8,0)--++(-12,0,0)--cycle;
    \draw (7,\boardY,7)--++(6,0,0)--++(0,5,0)--++(-6,0,0)--cycle;
    \foreach \pinX in {5,8,11,14}{
      \draw (\pinX,\boardY,6)--++(0,-1.4,0);
    }
  }
  % Simplified cable bundle retained as a few uncluttered paths.
  \draw[black!50] (10,44,7) .. controls (6,57,15) and (16,58,15) .. (23,49,14);
  \draw[black!50] (10,32,7) .. controls (0,36,11) and (-1,16,11) .. (4,3,8);
  % Moving bridge spans the storage field.
  \draw[fill=white] (20,2,10)--(25,2,10)--(25,60,10)--(20,60,10)--cycle;
  \draw (20,2,10)--(20,2,13)--(20,60,13)--(20,60,10);
  \draw (20,2,13)--(25,2,13)--(25,60,13)--(20,60,13);
  \draw[line width=1.5pt,black!55] (23,3,14)--(23,59,14);
  \draw[white,line width=.4pt] (23,3,14)--(23,59,14);
  \draw[line width=.8pt] (18,3,14)--(18,58,14);
  % Head carriage and tapered open receiving hopper.
  \draw[fill=white] (17,41,15)--(29,41,15)--(29,54,15)--(17,54,15)--cycle;
  \draw[fill=white,line width=.65pt]
    (17,41,27)--(31,41,27)--(31,55,27)--(17,55,27)--cycle;
  \draw[fill=white] (17,41,27)--(20,44,17)--(28,44,17)--(31,41,27)--cycle;
  \draw[fill=white] (31,41,27)--(28,44,17)--(28,52,17)--(31,55,27)--cycle;
  \draw (19,43,27)--(29,43,27)--(29,53,27)--(19,53,27)--cycle;
  \draw[black!45] (19,43,27)--(22,46,20)--(26,46,20)--(29,43,27);
  \draw[black!45] (26,46,20)--(26,50,20)--(29,53,27);
  \draw[fill=white] (20,44,16)--(29,44,16)--(29,52,16)--(20,52,16)--cycle;
  % Small servo body on the side of the head.
  \draw[fill=white] (14,43,17)--(18,43,17)--(18,49,17)--(14,49,17)--cycle;
  \draw (14,43,17)--(14,43,21)--(18,43,21)--(18,49,21)--(18,49,17);
  % Conveyor, on a separate stand to the right of the frame.
  \foreach \standX in {84,129}{
    \draw (\standX,32,0)--(\standX,32,12);
    \draw (\standX,54,0)--(\standX,54,12);
    \draw (\standX,32,0)--(\standX,54,0);
  }
  \draw[fill=white] (81,30,10)--(136,30,10)--(136,56,10)--(81,56,10)--cycle;
  \draw (81,30,10)--(81,30,14)--(136,30,14)--(136,30,10);
  \draw (136,30,14)--(136,56,14)--(136,56,10);
  \draw[fill=white,line width=.65pt] (81,31,14)--(136,31,14)--(136,55,14)--(81,55,14)--cycle;
  \draw[black!45] (83,31,14)--(83,55,14);
  \draw[black!45] (134,31,14)--(134,55,14);
  \draw[black!30] (108,31,14)--(108,55,14);
  % Arrow denotes the physical direction of feeding.
  \draw[line width=.75pt,->] (125,43,14.2)--(98,43,14.2);
  % Camera support and camera above the viewing zone.
  \draw[line width=1.3pt,black!60] (109,29,0)--(109,39,35);
  \draw[line width=.4pt,white] (109,29,0)--(109,39,35);
  \draw[fill=white] (102,37,35)--(115,37,35)--(115,43,35)--(102,43,35)--cycle;
  \draw[fill=white] (102,37,35)--(102,37,39)--(115,37,39)--(115,37,35)--cycle;
  \draw[fill=white] (115,37,35)--(115,37,39)--(115,43,39)--(115,43,35)--cycle;
  \draw (102,37,39)--(102,43,39)--(115,43,39);
  \draw[black!45,densely dashed] (108,40,34)--(99,37,14);
  \draw[black!45,densely dashed] (108,40,34)--(116,49,14);
  % Front frame rail, belt and motor housings.
  \draw[fill=white] (0,0,0)--(80,0,0)--(80,0,3)--(0,0,3)--cycle;
  \draw[line width=1.8pt,black!60] (2,2,8)--(79,2,8);
  \draw[line width=.45pt,white] (2,2,8)--(79,2,8);
  \draw[line width=.9pt] (3,6,9)--(79,6,9);
  \foreach \motorX in {0,19}{
    \draw[fill=white] (\motorX,-3,5)--++(8,0,0)--++(0,8,0)--++(-8,0,0)--cycle;
    \draw (\motorX,-3,5)--++(0,0,-5)--++(8,0,0)--++(0,8,0)--++(0,0,5);
    \draw (\motorX,-3,0)--++(8,0,0)--++(0,0,5);
    \draw (\motorX,-3,5)--++(0,8,0);
  }
  \draw[line width=1pt] (4,1,5)--(4,1,10);
  \draw[line width=1pt] (23,1,5)--(23,1,10);
  \draw[line width=.7pt] (4,1,10)--(77,1,10)--(77,57,10);
\end{scope}
% The leaders sit outside the geometry and are independent of its projection.
\draw[black!60] (143,33)--(153,27)--(153,19);
\draw[black!60] (119,65)--(127,79)--(137,79);
\draw[black!60] (46,64)--(37,77)--(24,77);
\draw[black!60] (76,27)--(89,14)--(104,14);
\draw[black!60] (27,44)--(15,52)--(6,52);
\foreach \labelX/\labelY/\labelNumber in {153/15/1,141/79/2,20/77/3,108/14/4,6/56/5}{
  \node[draw=black!80,fill=white,circle,inner sep=0pt,minimum size=7mm,
    font=\bfseries\fontsize{12bp}{12bp}\selectfont]
    at (\labelX,\labelY) {\labelNumber};
}
\end{tikzpicture}
};
  \putBlock{127}{124}{158}{\textbf{1} Конвейер\quad \textbf{2} Камера\quad
    \textbf{3} Головка XY\\
    \textbf{4} Девять ячеек\quad \textbf{5} Электроника}

  \draw[rulegray,line width=.4pt] (12,-148)--(285,-148);
  \putBlock{12}{152}{170}{\textbf{Команда}\enspace /\enspace группа Б01-506}
  \teamMember{12}{Малышев}{Павел Георгиевич}{Программная часть}
  \teamMember{70}{Войлов}{Антон Евгеньевич}{CAD и сборка}
  \teamMember{128}{Янковец}{Александр Сергеевич}{Сборка и электроника}
  \putBlock{12}{190}{180}{Ментор: Негодяев Андрей Сергеевич}

  \putBlock{199}{152}{38}{\textbf{GitHub}}
  \putBlock{247}{152}{38}{\textbf{Telegram}}
  \node[anchor=north west,inner sep=0pt] at (199,-161)
    {\qrcode[height=26mm,level=M]{https://github.com/blessedBA/sorting-warehouse}};
  \node[anchor=north west,inner sep=0pt] at (247,-161)
    {\qrcode[height=26mm,level=M]{https://t.me/sortirovachka}};
  \putBlock{195}{190}{45}{\href{https://github.com/blessedBA/sorting-warehouse}{Код проекта}}
  \putBlock{242}{190}{43}{\href{https://t.me/sortirovachka}{ПП Сортировщик}}
\end{scope}
\end{tikzpicture}
\newpage
\null
\begin{tikzpicture}[remember picture,overlay,x=1mm,y=1mm]
\begin{scope}[shift={(current page.north west)}]
  \pageFrame{Python + OpenCV\enspace /\enspace Arduino\enspace /\enspace HX711}{2}
  \draw[rulegray,line width=.4pt] (185,-25)--(185,-192);
  \putBlock{12}{25}{171}{\sectionLabel{02}{Технические характеристики}}
  \putBlock{12}{36}{52}{\subLabel{Сенсоры}
    Камера — наличие и цвет предмета.\\[2mm]
    Тензодатчик TAL 222 и АЦП HX711 — масса.}
  \putBlock{70}{36}{52}{\subLabel{Управление}
    Компьютер с Python/OpenCV — анализ и выбор ячейки.\\[2mm]
    Arduino Uno — управление механизмами.}
  \putBlock{128}{36}{51}{\subLabel{Актуаторы}
    Два шаговых двигателя NEMA17 — XY.\\[2mm]
    Мотор — конвейер.\\[2mm]
    Сервопривод SG90 — затвор головки.}
  \draw[rulegray,line width=.4pt] (12,-89)--(179,-89);
  \putBlock{12}{93}{167}{\textbf{9 ячеек / 3 весовых ряда}\par\vspace{2mm}
    0 < m < 5 г\quad /\quad 5 ≤ m < 12 г\quad /\quad m ≥ 12 г\par\vspace{2mm}
    Внутри ряда ячейка выбирается случайно. Цвет распознаётся, но на маршрут
    пока не влияет.}
  \putBlock{12}{124}{168}{\textbf{Функциональная схема}}
  \begin{scope}[shift={(12,-134)}]
    % Input inside the page TikZ picture; origin is its top-left corner.
\tikzset{
  function_box/.style={draw=ink,line width=.55pt,rounded corners=1mm,
    minimum height=11mm,align=center,inner sep=1mm,font=\bodyFont,fill=white},
  signal_arrow/.style={-{Stealth[length=2mm]},line width=.6pt},
  signal_link/.style={{Stealth[length=2mm]}-{Stealth[length=2mm]},line width=.6pt}}
\node[function_box,minimum width=32mm] (camera) at (17,-6) {Камера};
\node[function_box,minimum width=56mm] (computer) at (81,-6) {Компьютер / OpenCV};
\node[function_box,minimum width=36mm] (arduino) at (81,-28) {Arduino};
\node[function_box,minimum width=39mm] (weight_sensor) at (20,-28) {Тензодатчик\\+ HX711};
\node[function_box,minimum width=46mm] (servo) at (144,-28) {Сервопривод\\затвора};
\node[function_box,minimum width=48mm] (drivers) at (81,-50) {Драйверы};
\node[function_box,minimum width=46mm] (motors) at (144,-50) {Двигатели\\XY и ленты};
\draw[signal_arrow] (camera)--(computer);
\draw[signal_link] (computer)--node[right,font=\bodyFont]{USB} (arduino);
\draw[signal_arrow] (weight_sensor)--(arduino);
\draw[signal_arrow] (arduino)--(servo);
\draw[signal_arrow] (arduino)--(drivers);
\draw[signal_arrow] (drivers)--(motors);

  \end{scope}

  \putBlock{193}{25}{92}{\sectionLabel{03}{Как пользоваться}}
  \stepBlock{36}{1}{Подключить установку. Задать ноль координат, тарировать
    пустые весы и откалибровать их эталонным грузом.}
  \stepBlock{64}{2}{При запуске задать пороги массы и время подачи.
    Положить предмет; в консоли — \commandText{main start}.}
  \stepBlock{87}{3}{Следить за массой, цветом и ячейкой в журнале.
    Статус — \commandText{main status}, остановка — \commandText{main stop}.}

  \draw[rulegray,line width=.4pt] (193,-113)--(285,-113);
  \putBlock{193}{118}{92}{\sectionLabel{05}{Прогресс проекта}}
  \putBlock{193}{129}{92}{\textbf{Было заявлено}\par
    Склад на 9 ячеек, перемещение XY, сортировка по цвету и массе.}
  \putBlock{193}{150}{92}{\textbf{Сделано}\par
    Собрана механика; реализованы распознавание цвета, взвешивание,
    цикл по массе и журнал.}
  \putBlock{193}{176}{92}{\textbf{Осталось}\par
    Испытать полный цикл и точность; добавить выбор ячейки по цвету.}
\end{scope}
\end{tikzpicture}
\end{document}
